\section{Worum es geht}

A120 zeigt, dass die Fünffach-Drehung $R_5$ des Ikosaeders die drei
$Z_3$-Moden mit den Gewichten
\[
  p_0 = \tfrac29,\qquad p_1 = \tfrac{2+3\varphi}{9},\qquad p_2 = \tfrac{5-3\varphi}{9}
\]
umverteilt, und lässt eine Kante offen: wie diese Gewichte mit den
Massenverhältnissen der geladenen Leptonen zusammenhängen. Dieses Dokument
schließt die Kante, soweit der Korpus sie heute schließt. Es führt zwei Wege
aus, die beide reine Verhältnisse liefern, ohne Anker und ohne
$K_\text{frak}$: die Koide-Formel aus einem einzigen Matrixelement von $R_5$
und das $\varphi$-Skelett der Gewichte selbst. Auf dem zweiten Weg trifft
\[
  \frac{m_\tau}{m_e} = \Bigl(74+\frac1{45}\Bigr)\varphi^8 = 3477{,}469
\]
den Messwert auf $3\times10^{-5}$, innerhalb der Messunsicherheit. Ebenso
deutlich benennt das Dokument, was der Weg nicht leistet: Die Faktoren $74$
und $1/45$ sind beobachtet, nicht hergeleitet, und für $m_\mu/m_e$ gibt es
keinen $\varphi$-Weg auf diesem Niveau.

\textbf{[SETZUNG]} Vorausgesetzt ist wie in A120 die Einbettung der
$C_3$-Struktur in die Ikosaedergruppe $A_5$.

\section{Warum das Ikosaeder: die Drei und die Fünf}

\textbf{[B]} Die FFGFT ist eine $Z_3$-Theorie; der goldene Schnitt ist die
Zahl der Fünf. Die Drehgruppe $A_5$ des Ikosaeders ist die kleinste endliche
Gruppe, in der eine Dreier- und eine Fünferdrehung zusammen vorkommen, ohne
miteinander zu vertauschen. Aus diesem Nicht-Vertauschen entsteht die Mischung
der $Z_3$-Moden unter $R_5$, die A120 ausrechnet.

\textbf{[B]} Die beiden Symmetrien treten auf verschiedenen Ebenen ein. Die Drei
wirkt als Gitteroperation auf dem kompakten Träger: Die Trialität des
$D_4$-Gitters ist die $Z_3$ der Theorie. Eine Fünf kennt das Gitter nicht,
denn $|\mathrm{Aut}(D_4)| = 1152 = 2^7\cdot3^2$ und $5\nmid1152$; kein
Ikosaeder wirkt auf einem $D_4$-Gitter. Die Fünf tritt erst im ausgerollten
Raum $\mathbb{R}^3\times S^1$ (A090) ein, und zwar als Phase: Die
$72^\circ$-Drehung hat das charakteristische Polynom
$(x-1)\bigl(x^2-(\varphi-1)x+1\bigr)$ mit Spur $\varphi$. Modulo~3 liegt
$\varphi$ in $\mathrm{GF}(9)$ und ist dort primitiv ($\varphi^4=-1$,
Ordnung~8); die Eigenwerte $\zeta_5$ liegen in $\mathrm{GF}(81)$, weil
$5\mid80=|\mathrm{GF}(81)^*|$.

\textbf{[B]} Die Verteilung auf die beiden Ebenen ist erzwungen, nicht gewählt.
In Charakteristik~3 gilt $x^3-1=(x-1)^3$; in keinem $\mathrm{GF}(3^k)$ gibt es
eine primitive dritte Einheitswurzel. Die Drei kann deshalb keine Phase sein
und muss Gitter sein; die Fünf kann kein Gitter sein und muss Phase sein.
Drei und Fünf treffen sich nicht in der Geometrie, sondern im Körper
$\mathrm{GF}(81)$; die Erweiterung $\mathrm{GF}(9)\to\mathrm{GF}(81)$ ist die
Verdopplung $6=3\oplus3'$.

\section{Ein Matrixelement trägt die Koide-Formel}

\textbf{[B]} Das Gewicht $p_0=2/9$ und der Koide-Winkel $\theta=2/9$ (A110,
A120) sind derselbe Quotient $2/3^2$: im Zähler die Zahl der nicht-trivialen
$Z_3$-Moden, im Nenner das Quadrat der Gruppenordnung. Für das Matrixelement
$A=\langle v_0\,|\,R_5\,v_e\rangle$ gilt $|A|=\sqrt2/3$; der Betrag liefert
die Koide-Amplitude, das Betragsquadrat die Phase:
\[
  a_k = 1 + 3|A|\cos\bigl(|A|^2 + \tfrac{2\pi k}{3}\bigr),\qquad
  3|A| = \sqrt2,\qquad |A|^2 = \tfrac29 .
\]
Dieselbe Zahl erscheint einmal als Betrag und einmal als Betragsquadrat
desselben Matrixelements. Die Koide-Formel enthält damit keinen freien
Parameter mehr; die Massenamplituden folgen aus der $A_5$-Geometrie.

\textbf{[K]} Die Verhältnisse $(a_0/a_1)^2$ und $(a_2/a_1)^2$ ergeben
\begin{center}
\small
\begin{tabular}{@{}lllll@{}}
\toprule
Verhältnis & Formel & Wert & Messwert & Abweichung \\
\midrule
$m_\tau/m_e$ & $(a_0/a_1)^2$ & $3477{,}473$ & $3477{,}37\pm0{,}18$ & $+3{,}1\times10^{-5}$ ($0{,}6\sigma$) \\
$m_\mu/m_e$ & $(a_2/a_1)^2$ & $206{,}77032$ & $206{,}7682827(46)$ & $+9{,}8\times10^{-6}$ ($442\sigma$) \\
\bottomrule
\end{tabular}
\end{center}
Für $m_\tau/m_e$ liegt der Wert innerhalb der Messunsicherheit von $m_\tau$.
Für $m_\mu/m_e$ ist das exakte Paar $(\sqrt2,\,2/9)$ ausgeschlossen: Das
Verhältnis hängt bei fester Amplitude nur von $\theta$ ab, und die
Abweichung übersteigt die Unsicherheit um mehr als zwei Größenordnungen.
Das ist kein Einwand gegen die Herkunft von
$2/9$, sondern eine Grenze des exakten Paares.

\textbf{[B]} Weil $2/9\notin\pi\mathbb{Q}$, ist $\cos(2/9)$ transzendent. Exakte
algebraische Massenamplituden sind auf diesem Weg deshalb grundsätzlich
ausgeschlossen; jede algebraische Form kann den Koide-Wert nur annähern.

\section{Das $\varphi$-Skelett der Gewichte}

\textbf{[B]} Die beiden $\varphi$-haltigen Gewichte stehen in exakten Verhältnissen
zueinander und zum rationalen Gewicht:
\[
  \frac{p_1}{p_2} = \varphi^8,\qquad
  \frac{p_0}{p_2} = 2\varphi^4,\qquad
  \frac{p_1}{p_0} = \frac{\varphi^4}{2},\qquad
  3\sqrt{p_j}\in\bigl\{\sqrt2,\ \varphi^2,\ \varphi^{-2}\bigr\}.
\]
Das sind reine Zahlen aus der Geometrie, ohne Kosinus und ohne Skala. Die
Koide-Amplitude $\sqrt2=3\sqrt{p_0}$ steht in derselben Reihe wie $\varphi^2$
und $\varphi^{-2}$.

\section{Vom Skelett zu den Massenverhältnissen}

\textbf{[K]} Die Massenverhältnisse tragen dieses Skelett. Am genauesten trifft
\[
  \frac{m_\tau}{m_e} = \Bigl(74+\frac1{45}\Bigr)\varphi^8 = \frac{3331}{45}\,\varphi^8 = 3477{,}469 ,
\]
gegen den Messwert $3477{,}37\pm0{,}18$ (PDG~2024) also $+3{,}0\times10^{-5}$
oder $0{,}6\sigma$. Die Übersicht:
\begin{center}
\small
\begin{tabular}{@{}llll@{}}
\toprule
Verhältnis & Formel & Wert & gegen Messwert \\
\midrule
$m_\tau/m_e$ & $74\,\varphi^8$ & $3476{,}42$ & $-2{,}7\times10^{-4}$ ($-5{,}3\sigma$) \\
$m_\tau/m_e$ & $(74+\tfrac1{45})\,\varphi^8$ & $3477{,}469$ & $+3{,}0\times10^{-5}$ ($0{,}6\sigma$) \\
$m_\tau/m_e$ & $74\,\varphi^8\,(1+\tfrac{27}{12}\xi)$ & $3477{,}468$ & $+2{,}9\times10^{-5}$ ($0{,}6\sigma$) \\
$m_\tau/m_e$ & Koide, $\theta=2/9$ & $3477{,}473$ & $+3{,}1\times10^{-5}$ ($0{,}6\sigma$) \\
$m_\mu/m_e$ & $30\,\varphi^4 = 15\,p_0/p_2$ & $205{,}62$ & $-0{,}55\,\%$ \\
\bottomrule
\end{tabular}
\end{center}

\textbf{[K]} Auf dem $\varphi$-Weg ist $m_\tau/m_e$ damit so genau wie auf dem
Koide-Weg, als reines Verhältnis in natürlichen Einheiten. Der Weg ist
algebraisch: Die $\varphi$-Form nähert den transzendenten Koide-Wert an, und
die zweite Lage ($1/45$ bzw.\ $\tfrac{27}{12}\xi$) trägt den Rest. Gegen den
Koide-Wert selbst liegt $(74+\tfrac1{45})\varphi^8$ nur $1{,}2\times10^{-6}$
daneben.

\textbf{[S]} Die Faktoren sind beobachtet, nicht hergeleitet. $74=2\cdot37$ ist
dieselbe Zahl wie im Zähler der Rotationszahl $74/75=1-100\xi$ der Rekursion
(A050) und im Wert von $K_\text{frak}$ (A040). Die zweite Lage hat dieselbe
Form wie die Galois-Form $1/\alpha=3700/27=137+\tfrac1{27}$ (A130): ganze Zahl
plus Stammbruch. Die~37 steht dabei nur in Zählerteilen ($74=2\cdot37$,
$3700=100\cdot37$), in keinem Nenner ($45=3^2\cdot5$, $27=3^3$); eine
gemeinsame algebraische Quelle ist damit nicht belegt.

\textbf{[B]} Aus dem Ikosaeder kommt die~37 nachweislich nicht. Im Phasenkanal
ist sie die Einheitswurzel $\zeta_{37}$, und diese lebt in
$\mathrm{GF}(3^{18})$, weil die Ordnung von~3 modulo~37 gleich~18 ist. Die
ikosaedrische Phase $\zeta_5$ lebt in $\mathrm{GF}(81)=\mathrm{GF}(3^4)$, und
$4\nmid18$: $\mathrm{GF}(81)$ ist kein Teilkörper von $\mathrm{GF}(3^{18})$.
Gemeinsam haben beide Körper nur $\mathrm{GF}(9)$, also $\varphi$.

\section{Zwei Lesarten derselben Zahlen}

\textbf{[B]} Die Leiter aus A100 und der Koide-Weg geben dieselben Zahlen, aber
sie lassen sich in zwei Richtungen lesen. In Lesart~A ist die
Leiter $r_k\,\xi^{p_k}$ fundamental, und der Schritt vom Prozent- auf das
$10^{-5}$-Niveau ist eine gemessene Korrektur. In Lesart~B ist die
Koide-Formel aus dem Ikosaeder fundamental; dann ist die Leiter eine rationale
Näherung an die transzendenten Amplituden $a_k^2$, und die generationsabhängigen
Reste sind der Abstand zwischen Kosinus und Potenzgesetz, ausgerechnet und
nicht gemessen \textbf{[K]}. Sie fallen für die drei Leptonen verschieden aus,
weil die drei Kosinuswerte verschieden sind. Welche Lesart die tragende ist,
entscheidet die Rechnung nicht \textbf{[S]}.

\section{Was der Weg nicht leistet}

\textbf{[S]} Für $m_\mu/m_e$ gibt es keinen $\varphi$-Weg auf dem Niveau von
$m_\tau/m_e$: $30\varphi^4$ liegt $0{,}55\,\%$ daneben, die Leiter aus A100
$0{,}52\,\%$, und das exakte Koide-Paar ist ausgeschlossen (oben).

\textbf{[K]} Das $\varphi$-Skelett ist nicht die frühere $\varphi$-Massenleiter
mit $m\propto\varphi^{\text{Generation}}$ und $m_\mu/m_e\approx\varphi^{10}$.
Diese ist widerlegt: $\varphi^{10}=122{,}99$, nicht $206{,}77$.

\section{Prüfskript}

\begin{itemize}
  \item Gewichte $p_j$ aus $R_5$ (exakt gegen $2/9$, $(2+3\varphi)/9$,
    $(5-3\varphi)/9$), $|A|=\sqrt2/3$, das $\varphi$-Skelett, die
    Koide-Verhältnisse gegen PDG~2024 und CODATA, die Tabelle der
    $\varphi$-Formen, die Körperaussagen ($5\nmid1152$, $\varphi^4=-1$ in
    $\mathrm{GF}(9)$, Ordnung von~3 modulo~37, $4\nmid18$) und die Abgrenzung
    zu $\varphi^{10}$:
    \texttt{python/A\_Serie\_Skripte/a125\_phi\_skelett.py}
\end{itemize}

\section{Zeichenerklärung}

Symbole, die in der gesamten A-Serie verwendet werden, sind in A010 erklärt.
Spezifisch für dieses Dokument:

\begin{center}
\renewcommand{\arraystretch}{1.3}
{\small\begin{tabular}{lp{9cm}}
\toprule
Symbol & Bedeutung \\
\midrule
$\varphi$ & goldener Schnitt $(1+\sqrt5)/2$ \\
$A_5$, $R_5$ & Drehgruppe des Ikosaeders; Fünffach-Drehung um $72^\circ$ \\
$v_0$, $v_e$ & triviale und nicht-triviale (Elektron-)Mode des $Z_3$-Zirkulanten \\
$p_j$ & Umverteilungsgewichte $|\langle v_j\,|\,R_5\,v_e\rangle|^2$ \\
$A$ & Matrixelement $\langle v_0\,|\,R_5\,v_e\rangle$, $|A|=\sqrt2/3$ \\
$a_k$ & Koide-Amplituden $1+\sqrt2\cos(\theta+2\pi k/3)$ \\
$\zeta_n$ & primitive $n$-te Einheitswurzel \\
\bottomrule
\end{tabular}}
\renewcommand{\arraystretch}{1.0}
\end{center}

\section{Was offen bleibt}

\textbf{Erstens die Faktoren $74$ und $1/45$.} Sie stammen nicht aus den
Gewichten selbst und nachweislich nicht aus dem Ikosaeder. Ob sie aus der
Galois-Hierarchie $\mathrm{GF}(3^n)$ folgen, in der auch $K_\text{frak}$ und
$3700/27$ stehen, ist offen.

\textbf{Zweitens die gemeinsame Wurzel.} Das Ikosaeder liefert Phasen in
$\mathrm{GF}(81)$, die Zahlentheorie liefert Ordnungen in
$\mathrm{GF}(3^{18})$; beide treffen die Leptonmassen, aber sie berühren
sich nur in $\varphi$. Ob es eine Struktur gibt, in der beide dieselbe Wurzel
haben, oder ob es zwei unabhängige Zugänge sind, ist die offene Kante.

\textbf{Drittens $m_\mu/m_e$.} Für dieses Verhältnis fehlt ein algebraischer
Weg auf dem Niveau von $10^{-5}$.

\textbf{Viertens die Einbettung $C_3\subset A_5$.} Sie ist wie in A120
vorausgesetzt.

\section{Ersetzt}

Dieses Dokument nimmt auf: Dok.~367 (ein Matrixelement trägt Koide),
Dok.~368 ($\varphi$-Skelett und Galois-Faktoren), Dok.~369 (Wege~2 und~4 der
Bestandsaufnahme), Dok.~370 (warum das Ikosaeder), Dok.~364 (soweit $A_5$ und
$D_4$ betreffend). Eingearbeitet: R121. Abgegrenzt: Dok.~172, 281 (widerlegte
$\varphi$-Massenleiter).

\section{Verwendung von KI-Werkzeugen}

Dieses Dokument ist unter Verwendung KI-gestützter Werkzeuge entstanden. Die
Fragestellungen, die Setzungen und alle inhaltlichen Entscheidungen stammen vom
Autor; die Werkzeuge formulieren aus, rechnen nach und erstellen Prüfskripte.
Die Verantwortung für Inhalt und Fehler liegt vollständig beim Autor. Der
ausführliche Wortlaut steht in A010.
